% --- PROBLEMA 4: BASES INTRÍNSECAS ---
\subsection{Problema 4: Curvas Parametrizadas (Base de Serret-Frenet)}
Para la siguiente curva en el espacio, obtenga los vectores unitarios Tangente ($\hat{T}$), Normal ($\hat{N}$) y Binormal ($\hat{B}$):

$$\vec{r}(s) = \left( a \cos\left(\frac{s}{c}\right), a \sin\left(\frac{s}{c}\right), \frac{b s}{c} \right)$$ 
donde $c = \sqrt{a^2 + b^2}$ y $s$ es el parámetro de longitud de arco.

\begin{enumerate}
    \item Obtenga $\hat{T} = \frac{d\vec{r}}{ds}$.
    \item Obtenga $\hat{N} = \frac{d\hat{T}/ds}{\|d\hat{T}/ds\|}$.
    \item Obtenga $\hat{B} = \hat{T} \times \hat{N}$.
    \item Calcule la curvatura $\kappa = \|\frac{d\hat{T}}{ds}\|$.
\end{enumerate}